\subsection{Sprint 4}

\subsubsection*{Planejamento e Distribuição das Atividades}

Na definição das tarefas da Sprint 4, a Squad ficou responsável pelo desenvolvimento
da \textbf{Task 16 -- Gestão de Vagas do Usuário Empresa}. A distribuição das
atividades ocorreu da seguinte forma: o desenvolvimento do \textit{frontend} ficou
sob minha responsabilidade, enquanto os integrantes Mayra e Vinícius ficaram
responsáveis pelas atividades de \textit{backend}. Adicionalmente, foi atribuída a
mim uma atividade de débito técnico referente à \textbf{Task 5.0 -- Cursos Internos},
relacionada ao cadastro de cursos internos para alunos.

\subsubsection*{Débito Técnico -- Task 5.0: Cursos Internos}

Após a definição das tarefas, foi realizada a atualização do ambiente de
desenvolvimento com o objetivo de iniciar a implementação da Task 5.0. Durante esse
processo, as imagens antigas do Docker foram removidas e o ambiente foi reconstruído
integralmente, seguido da execução de testes para validação do funcionamento da
aplicação.

Ao iniciar a análise da funcionalidade, observou-se que os requisitos descritos na
Task 5.0 já estavam implementados no ambiente de desenvolvimento. Foi realizada uma
verificação detalhada dos fluxos da aplicação, confirmando que os critérios de
aceitação previstos estavam atendidos. Dessa forma, concluiu-se que a funcionalidade
já havia sido desenvolvida anteriormente.

Durante a revisão da implementação, foi identificado um problema no fluxo de inscrição
em cursos externos. Embora a informação referente ao link de inscrição estivesse
corretamente armazenada no banco de dados e retornada pela \ac{api}, verificou-se que
esse campo não estava sendo utilizado na interface \textit{frontend}. Após a
confirmação da situação por meio de consultas realizadas no Postman, foi encaminhada
uma comunicação ao integrante responsável da \ac{ages} IV relatando que a
funcionalidade principal já se encontrava implementada, restando apenas a correção
relacionada à exibição do link de inscrição.

Posteriormente, foi constatado que o desenvolvimento original da funcionalidade havia
sido realizado por um integrante da \ac{ages} III. Diante disso, a inconsistência
identificada foi repassada a ele, permitindo que a correção fosse realizada em conjunto
durante a aula. Após a implementação da exibição do link de inscrição, a alteração foi
integrada ao código do projeto.

\subsubsection*{Task 16.1 -- Desenvolvimento do Frontend de Gestão de Vagas}

No dia 09/06, teve início o desenvolvimento da Task 16.1 -- Gestão de Vagas para
Usuário Empresa, referente ao \textit{frontend} da funcionalidade. Durante a análise
inicial, foi identificada a existência de uma \textit{branch} no repositório contendo
uma implementação parcial da tela. Após verificação, constatou-se que a
\textit{branch} havia sido iniciada anteriormente pela integrante Ana, da \ac{ages}
III, porém ainda não havia sido submetida para aprovação por meio de
\textit{Pull Request}.

Considerando a existência dessa implementação preliminar, foi realizada uma análise
comparativa entre os requisitos descritos no ClickUp, o protótipo disponibilizado no
Figma e os padrões adotados em outras telas do sistema. Essa análise permitiu
identificar quais elementos poderiam ser reaproveitados e quais necessitavam de
adequação para atender aos critérios estabelecidos.

As principais alterações realizadas envolveram a refatoração da tela de Gestão de
Vagas da Empresa, a padronização dos filtros de pesquisa, a correção da paginação,
a adequação dos ícones utilizados e a reorganização dos componentes visuais para
alinhamento com o \textit{design} definido no Figma. Também foram realizadas
verificações de conformidade com os padrões arquiteturais e visuais do projeto,
seguidas da execução de testes manuais para validação do funcionamento da
funcionalidade.

Após a conclusão das atividades, foi aberto o \textit{Pull Request} referente às
alterações desenvolvidas. Em seguida, o integrante HArthur, da \ac{ages} IV, informou
que realizaria uma refatoração de componentes compartilhados, uma vez que havia
desenvolvido novos componentes que já estavam sendo utilizados em outras páginas do
sistema. O acompanhamento dessa atividade foi realizado de forma conjunta e, após
aproximadamente trinta minutos de ajustes e validações, a \textit{branch} foi integrada
à \textit{branch} Development do projeto.

\subsubsection*{Task 16.2 -- Suporte ao Desenvolvimento do Backend}

No dia 11/06, foi realizado contato com os integrantes da Squad para verificar o
andamento das atividades de \textit{backend} relacionadas à Task 16.2 e suas
respectivas subtarefas. Durante a conversa, foi informado que um dos integrantes não
conseguiria concluir o desenvolvimento dentro do prazo previsto e, considerando a
proximidade da data de entrega, foi solicitada minha colaboração para auxiliar na
implementação.

Ao iniciar a análise da funcionalidade, em um primeiro momento não foram identificadas
pendências evidentes. Entretanto, considerando que se tratava de uma funcionalidade
associada a uma nova tela e que ainda não havia sido desenvolvida pelos integrantes
responsáveis, foi realizada uma verificação mais aprofundada, comparando os requisitos
descritos no ClickUp com a implementação existente.

Durante essa análise, foi identificada a ausência do campo \textit{Work Type},
responsável por indicar a modalidade da vaga, podendo assumir os valores
\textit{presencial}, \textit{remoto} ou \textit{híbrido}. Por meio da análise da
estrutura do banco de dados PostgreSQL, foram identificados os campos correspondentes,
permitindo a implementação da funcionalidade necessária para disponibilizar essa
informação corretamente.

Além disso, foi constatado que a funcionalidade de criação de vagas não estava
retornando adequadamente as habilidades (\textit{Skills}) associadas à vaga. Após a
análise da lógica implementada, verificou-se que a consulta retornava apenas os
identificadores (UUIDs) provenientes da tabela de relacionamento \textit{SkillJob},
sem fornecer as informações necessárias para exibição na aplicação. Dessa forma, foram
realizados os ajustes necessários para retornar corretamente os dados relacionados às
habilidades, tornando a funcionalidade compatível com os requisitos definidos.

Após a conclusão das alterações, foi aberto um \textit{Pull Request} e executadas as
validações automáticas disponibilizadas pelo projeto, incluindo as verificações
realizadas pelo GitHub Actions e GitHub Copilot, sem que fossem identificados erros de
implementação.

Posteriormente, durante uma conversa com o integrante responsável da \ac{ages} IV,
foi informado que as funcionalidades já haviam sido implementadas anteriormente e que
a versão disponível no ambiente Development correspondia a uma atualização recente
realizada por ele. Diante dessa informação, foram apresentados os pontos identificados
durante a análise, especificamente as possíveis inconsistências relacionadas ao campo
\textit{Work Type} e ao retorno das \textit{Skills}. Após nova validação da equipe,
foi confirmado que as pendências identificadas já haviam sido solucionadas, resultando
na conclusão das atividades sob minha responsabilidade.

\subsubsection*{Semana de Freeze -- Suporte a Débito Técnico}

Durante a semana de Freeze do projeto, fui procurado por integrantes da Squad em razão
de dificuldades enfrentadas por uma colega da \ac{ages} I na execução de uma tarefa de
débito técnico. Com o objetivo de auxiliar na resolução do problema, foi realizada a
atualização do ambiente de desenvolvimento e uma análise da funcionalidade em questão.

Após a investigação do comportamento apresentado, foi identificada uma possível causa
para o problema e elaborada uma sugestão de solução. As orientações técnicas sobre os
ajustes que poderiam ser realizados foram repassadas à colega, que aplicou a abordagem
sugerida com sucesso, permitindo a conclusão da atividade.

\subsubsection*{Avaliação da Sprint}

A Sprint 4 evidenciou a importância da comunicação contínua entre os integrantes da
equipe e entre as diferentes turmas do projeto. Em mais de uma situação, a troca de
informações com colegas de \ac{ages} III e \ac{ages} IV foi determinante para evitar
retrabalho e garantir a consistência das entregas.

A experiência de atuar tanto no \textit{frontend} quanto no \textit{backend} ao longo
da sprint proporcionou uma visão mais abrangente da funcionalidade desenvolvida, além
de reforçar a capacidade de adaptação diante de necessidades imprevistas da equipe. A
identificação proativa de inconsistências -- mesmo em atividades que não eram de
responsabilidade direta -- contribuiu para a qualidade geral do produto entregue.

Por fim, o suporte prestado à colega durante a semana de Freeze reforçou a dinâmica
colaborativa da equipe, demonstrando que a solidariedade entre os integrantes, mesmo
em períodos de menor atividade formal, é um fator relevante para o sucesso do projeto
como um todo.
